\mode*
\subsection{Problem solving and program dynamics}
\label{problem-solving}

\subsubsection{Role in the syllabus}

Problem solving is, in a sense, the apex skill of CS1: taking a problem
description, decomposing it into sub-problems, implementing each, and assembling
the pieces into a working program. It is also where many students fall short.
The ITiCSE working group led by \textcite{Lister2004Tracing} established that, at
the end of an introductory course, a large fraction of students still cannot
program. A natural explanation is a missing problem-solving ability; but the
working group's evidence points one level lower, to the more basic skills that
problem solving rests on:
\blockcquote{Lister2004Tracing}{many students do not know how to program at the
conclusion of their introductory courses. A popular explanation for this
incapacity is lack of the ability to problem-solve [\ldots] An alternative
explanation is that many students have a fragile grasp of both basic programming
principles and the ability to systematically carry out routine programming
tasks, such as tracing [\ldots] suggesting that these skills are a prerequisite
for problem-solving.}
Follow-up work sharpened this into a hierarchy: student performance on tracing
correlates with performance on code writing, and a path analysis
\textcquote{Lopez2008Hierarchy}{suggests the possibility of a hierarchy of
programming related tasks}, with knowledge of programming constructs at the
bottom, tracing at an intermediate level, and writing at the top
\autocite{Lopez2008Hierarchy}.

This matters for our analysis: if tracing and a firm grasp of individual
constructs are prerequisites for problem solving, then teaching against the
construct-level misconceptions catalogued above is also groundwork for problem
solving. What remains for this section is the misconception that blocks the
intermediate level itself.

\subsubsection{The program is not its text}

The misconception underneath poor tracing is a static reading of the program:
\textcite[p.~235]{Sorva2012Dissertation} describes
\textcquote[p.~235]{Sorva2012Dissertation}{some learners' failure to see that a
program has a dynamic existence in addition to its static aspect}.
Read that way, the text order \emph{is} the execution order.
\textcite{Sleeman1984} found novices who held exactly this rule:
\textcquote[p.~19]{Sleeman1984}{All statements including those in
procedure bodies were executed in the order they appeared}, and difficulties
with the sequentiality of statements recur across the literature
\autocite[Table~A.1, p.~360]{Sorva2012Dissertation}.
The static reading has a classic sibling, which \textcite{Pea1986Superbug}
named the \emph{parallelism} bug:
\textcquote[p.~27]{Pea1986Superbug}{the assumption that different lines in a
program can be active or somehow known by the computer at the same time, or
in parallel}. Pea traces his three bug classes --- parallelism,
intentionality, egocentrism --- to a common root, a
\textcquote[p.~32]{Pea1986Superbug}{\enquote{superbug,} the default strategy that
there is a hidden mind somewhere in the programming language that has
intelligent interpretive powers}. Text-order execution and parallel
execution miss the same critical aspect from opposite sides: exactly one
line is active at a time, and which one is decided by control flow alone.
Students struggle to follow code step by step and to reason about the order in
which a program will execute \autocite[p.~20]{Sleeman1984}, a difficulty also
documented by \textcite{Kurvinen2016}. Both feed directly into the
problem-solving gap that \textcite{Lister2004Tracing} identified, since a
student who cannot trace a short program reliably has little hope of assembling
a larger one.

\mode<presentation>{%
\begin{frame}
  \begin{block}{Misconception\autocite{Sleeman1984,Pea1986Superbug,Sorva2012Dissertation}}
    \begin{itemize}
      \item A program is read as static text: statements execute in the
        order they are written --- including the bodies of functions, where
        they stand --- or are all \enquote{active} at once.
    \end{itemize}
  \end{block}
\end{frame}%
}

The object of learning is the \emph{execution order} of a program, and the
critical aspect is that control flow --- not textual position --- decides
when, and whether, a line runs.
The contrast keeps the program's behaviour invariant and varies only the
textual arrangement: both programs below print
\mintinline{python}{a}, \mintinline{python}{b}, \mintinline{python}{c}, in
that order.
A student reading text order as execution order predicts the same for the
left program but \mintinline{python}{b}, \mintinline{python}{a},
\mintinline{python}{c} for the right one --- the line printing
\mintinline{python}{b} stands first in the text, yet runs second, because a
definition is not an execution site.

\hfill
\begin{minipage}[t]{0.45\columnwidth}
  \begin{minted}{python}
    print("a")
    print("b")
    print("c")
  \end{minted}
\end{minipage}
\hfill
\begin{minipage}[t]{0.45\columnwidth}
  \begin{minted}[highlightlines={1,2}]{python}
    def middle():
        print("b")
    print("a")
    middle()
    print("c")
  \end{minted}
\end{minipage}

To generalise, we keep the critical value invariant --- the execution order
is derived from control flow, not read off the page --- and vary the
construct that makes them diverge, one per analysis in the preceding
sections: a false \mintinline{python}{if} \emph{skips} its suite, a loop
\emph{revisits} its body, a call \emph{jumps} into a body defined elsewhere
and returns. Tracing each with a finger on the currently-executing line
makes the deviation from text order the invariant experience across all
three.

For the fusion, the student must hold these together in one program --- and
there is direct evidence that this integration is a separate difficulty, not
a free consequence of mastering the parts.
\textcite{Huang2024Composition} found in novice code tracing
\textcquote{Huang2024Composition}{a robust composition effect across problems
and students}: problems requiring several component skills at once are harder
than the matched single-skill problems, with
\textcquote{Huang2024Composition}{conceptual misunderstandings of how to
integrate component skills} as one identified source of errors.
In variation-theory terms, fusion is where the integration is taught rather
than assumed. The function from the return-statement analysis in
\cref{misconceptions} serves: a definition (jumped into), a loop (revisited),
a conditional (skipped or not) and an early return, traced line by line.

\begin{minted}{python}
    def first_negative(values):
        for value in values:
            if value < 0:
                return value
        return None

    print(first_negative([4, -1, -5]))
\end{minted}

Diagnosing tracing misconceptions is itself an active line of work:
\textcite{Bastian2026AdaptiveTracing} construct an adaptive test to
\textcquote{Bastian2026AdaptiveTracing}{identify these misconceptions}
related to control flow. Our diagnostic instruments
(\cref{app:diagnostics}) take the same stance with fixed items: each
distractor encodes one catalogued misconception.

\subsubsection{Difficulties that can occur}

Beyond the misconception above, two difficulties recur. The first is
mathematical content.
\textcite{KumarVeerasamy2016} found that students did markedly worse on
math-related tasks than on others:
\blockcquote{KumarVeerasamy2016}{This study analysis also explored that novices
of programming struggled in writing code for math related questions 6 and 7
\textelp{}. Nearly 66\% of students did not do well in the mathematical problems
based questions though explained and allowed to surf the Internet to seek for
more details during the exam hours. A Neo-Piagetian theory of cognitive
development stated that students who are at the concrete operational stage
struggle to write large programs with partial specifications, although they can
write small programs from well-defined specifications \textelp{}.}

The second is the integration difficulty already discussed: even students who
handle each construct separately may fail when a problem composes them
\autocite{Huang2024Composition}. Neither is a misconception in the
catalogue's sense --- there is no false belief to contrast away --- but both
constrain how far the patterns of variation can carry a student, and the
composition effect in particular motivates never stopping a teaching design
before its fusion step.
